%=========================================================================
%  THEME 4 — "RUECKGRAT"
%
%  Eine senkrechte Linie läuft als Rückgrat durch jede Seite. Links davon
%  stehen rechtsbündig Daten, Abschnittstitel und Schlüssel, rechts der
%  Text. Grüne Knoten auf der Linie markieren die Abschnitte. Gesetzt in
%  JetBrains Mono, einer Schrift mit fester Zeichenbreite.
%
%  Gestaltung: Markus Neusinger (github.com/MarkusNeusinger, anyplot.ai),
%  ursprünglich als eigene Dokumentklasse rueckgrat.cls. Hier als Theme
%  umgesetzt, das dieselben Inhaltsdateien liest wie die anderen Layouts.
%
%  Die linke Spalte fasst rund 22 Zeichen je Zeile. Eine längere
%  E-Mail-Adresse ragt über die Linie: dann \rgcolL unten vergrössern.
%=========================================================================

% Masse des Rasters: linker Seitenrand, linke Spalte, Abstand zur Textspalte
\newlength{\rgmarginleft}\setlength{\rgmarginleft}{15mm}
\newlength{\rgcolL}\setlength{\rgcolL}{37mm}
\newlength{\rgcolG}\setlength{\rgcolG}{6mm}
\newlength{\rgpull}\setlength{\rgpull}{\dimexpr\rgcolL+\rgcolG\relax}
\newlength{\rgphotow}\setlength{\rgphotow}{32mm}

% Der Satzspiegel ist die Textspalte. Die linke Spalte steht im Seitenrand.
\usepackage[left=58mm, right=18mm, top=18mm, bottom=20mm,
            headheight=12.5pt, headsep=5mm, footskip=9mm]{geometry}

%=========================================================================
%  GEMEINSAME BASIS — Pakete, Icons, Hilfsmakros.
%  Enthält bewusst KEINE Gestaltung. Alles Gestalterische (Farben, Schrift,
%  Raster, Abstände) liegt im jeweiligen Theme unter templates/themes/.
%=========================================================================

\usepackage[british, ngerman]{babel}
\usepackage{etoolbox}
\usepackage{xcolor}
\usepackage{fontspec}
\usepackage[final]{microtype}

%-------------------------------------------------------------------------
%  Echte Leerzeichen im PDF. TeX setzt Wortzwischenräume sonst
%  nur als Abstand. Ältere PDF-Leser (z.B. pdf.js 2.x im PDF-Viewer von
%  VS Code) und manche Bewerbungsportale lesen den Text dann ohne Leerzeichen
%  aus ("Bei500TaktenproMinute"). Mit 2 schreibt XeTeX die Leerzeichen als
%  Glyphen mit. Seitenzahlen und Blocksatz bleiben gleich, einzelne Absätze
%  brechen etwas anders um (weniger Trennungen).
%  Gemessen mit pdf.js 2.6, Anteil einzeln lesbarer Wörter: Briefe vorher
%  3-5 %, nachher 97-99 %. Im Lebenslauf stieg der Anteil ebenfalls deutlich.
%  Rest: Zeilen in der Kompetenzentabelle und einzelne Zeilen mit Bindestrich
%  verlieren die Zeichen weiterhin (XeTeX-intern, nicht weiter eingegrenzt).
%  pdftotext, pdf.js 4 und gängige Parser setzen dort ohnehin Leerzeichen aus
%  den Abständen, ohne Doppelungen.
%  ACHTUNG, XeTeX 0.999993 stürzt damit ab (Segmentation fault), wenn
%    - irgendwo \spaceskip ungleich null ist (deshalb der Umbau bei ragged2e
%      unten) oder
%    - \RaggedRight/\bodyflow mitten in einem schon begonnenen Absatz
%      eingeschaltet wird ("Text {\bodyflow ...}"). Immer am Absatzanfang.
%  Und: Farbwechsel mitten im Fliesstext in derselben Schrift verrutschen
%  (das farbige Wort bleibt schwarz, der Text danach wird farbig). Abhilfe:
%  \kern0pt an beiden Enden des farbigen Stücks, z.B.
%  {\color{red}\kern0pt Text\kern0pt} (so \platz in feedback.tex.
%  Farbe zusammen mit Schriftwechsel (Ortsangaben, Symbole) ist nicht betroffen.
%  Stürzt ein Build nach einer Vorlagenänderung ab, zuerst diese beiden Punkte
%  prüfen. Notausgang: hier \XeTeXinterwordspaceshaping=0 setzen.
%-------------------------------------------------------------------------
\XeTeXinterwordspaceshaping=2
% Als Zeichen geschrieben wird nur ein normaler Wortabstand. Nach Grossbuch-
% staben (sfcode 999) und Satzzeichen (ohne \frenchspacing) setzt TeX einen
% leicht abweichenden Abstand und XeTeX lässt das Zeichen weg ("SAPHANACloud").
% Deshalb alle auf 1000 -- im Deutschen ohnehin üblich, im Satzbild unsichtbar.
\count255=`A
\loop\sfcode\count255=1000 \ifnum\count255<`Z \advance\count255 by 1 \repeat
\sfcode`Ä=1000 \sfcode`Ö=1000 \sfcode`Ü=1000
\AtBeginDocument{\frenchspacing}
% Unsichtbares Leerzeichen ohne Breite: dort, wo im Satz kein Wortabstand
% steht und Leser trotzdem einen brauchen -- zwischen zwei Spalten
% ("EntwicklungPython"), nach Aufzählungszeichen und am Schriftwechsel nach
% den Leads ("Prüfstandsdatenplattform:Von"). Ändert das Layout nicht.
\newcommand{\cvrealspace}{\ifnum\XeTeXcharglyph"20>0
    \rlap{\XeTeXglyph\XeTeXcharglyph"20\relax}\fi}
\usepackage{tikz}
\usetikzlibrary{calc}
\usepackage{graphicx}
\usepackage{enumitem}
\usepackage{fancyhdr}
\usepackage{tabularx}
\usepackage{array}
\usepackage[fixed]{fontawesome5}
\usepackage{paracol}
\usepackage{setspace}
\usepackage{ragged2e}
% ragged2e hält im Flattersatz die Wortabstände starr, indem es bei jedem
% Schriftwechsel \spaceskip auf den natürlichen Wortabstand setzt. Zusammen
% mit echten Leerzeichen (\XeTeXinterwordspaceshaping=2, oben) stürzt XeTeX
% 0.999993 dabei ab. Gleicher Effekt ohne \spaceskip:
% Dehnung und Stauchung des Wortabstands in der Schrift selbst auf null setzen
% (\fontdimen3/4). Der Wortabstand ist damit derselbe wie bei ragged2e.
% Nebenwirkung: Die Änderung gilt für die ganze Satzlauf-Instanz der Schrift.
% Unschädlich, weil die Briefe im Blocksatz vor dem Lebenslauf gesetzt werden
% bzw. in eigenen Läufen (make dossier).
\makeatletter
\renewcommand{\@raggedtwoe@everyselectfont}{%
    \spaceskip\z@
    \if@raggedtwoe@spaceskip
        \fontdimen3\font=\z@
        \fontdimen4\font=\z@
    \fi}
\makeatother
\usepackage{needspace}
\usepackage{accsupp}

% Textschicht für Bewerbungsportale. SmartRecruiters & Co.
% lesen den Lebenslauf aus dem PDF aus. Ohne diese Zeile kamen gesperrte
% Überschriften als "BER UFS E R FAH R UNG" und jeder Bindestrich als
% geschützter Bindestrich (U+2011) an. \XeTeXgenerateactualtext legt zu
% jedem Textstück den Originaltext ab; am Druckbild ändert sich nichts.
\ifdefined\XeTeXgenerateactualtext\XeTeXgenerateactualtext=1\fi


%-------------------------------------------------------------------------
%  Sprache
%
%  \cvlanguage ist ein babel-Sprachname und wird VOR dem Theme gesetzt --
%  in applications/<app>/meta.tex, denn die Sprache gehoert zur Ziel-
%  bewerbung und nicht zur Gestaltung. Voreinstellung ngerman: alle
%  bestehenden Bewerbungen bauen unveraendert, ohne dass ihre meta.tex
%  etwas davon wissen muss. Die babel-Optionen oben laden alle drei
%  Sprachen; die letzte ist die Hauptsprache, \selectlanguage schaltet um.
%  Fuer Englisch "british" verwenden, nicht "english": babel warnt vor der
%  mehrdeutigen Option, und british liefert das europaeische Datumsformat.
%-------------------------------------------------------------------------

\providecommand{\cvlanguage}{ngerman}

%-------------------------------------------------------------------------
%  Beschriftungen der Abschnitte -- Voreinstellung Deutsch.
%  Eine uebersetzte Inhaltsdatei (templates/themes/cv-content-en.tex)
%  setzt sie per \renewcommand um; kein Theme enthaelt mehr festen Text.
%-------------------------------------------------------------------------

%  Briefdatum. Voreinstellung ist babels \today. Die englische Inhalts-
%  datei setzt es um: babels british schreibt "16th September 2026", und
%  die Ordnungszahl ist im Geschaeftsbrief unueblich.
\providecommand{\lettertoday}{\today}

\providecommand{\LblSkills}{Kompetenzen}
\providecommand{\LblExperience}{Berufserfahrung}
\providecommand{\LblEducation}{Ausbildung}
\providecommand{\LblFurther}{Weiterbildung}
\providecommand{\LblLanguages}{Sprachen}
\providecommand{\LblContact}{Kontakt}
\providecommand{\LblProfile}{Profil}
\providecommand{\LblCV}{Lebenslauf}
\providecommand{\LblPage}{Seite}

\setlength{\parindent}{0pt}
\setlength{\parskip}{0pt}
\emergencystretch=2em

%-------------------------------------------------------------------------
%  Typografische Hygiene
%-------------------------------------------------------------------------

% Trennung in Überschriften/Titeln unterdrücken — verhindert Unfälle wie
% "MUSTERMANN TECH-|NIK AG" quer über den Zeilenumbruch.
\newcommand{\nohyph}{\hyphenpenalty=10000\exhyphenpenalty=10000\hbadness=10000}

% Eigennamen, die nie getrennt werden sollen. \hyphenation traegt die
% Ausnahmen nur in die gerade aktive Sprache ein -- deshalb als Makro und
% nach \selectlanguage noch einmal aufgerufen.
\newcommand{\cvhyphenexceptions}{%
    % Eigene Firmen-, Orts- und Produktnamen hier ergänzen. Ohne Bindestrich
    % = nie trennen, mit Bindestrich = nur an dieser Fuge trennen.
    \hyphenation{ANSYS FastAPI Team-center Work-bench}}
\cvhyphenexceptions
% \frenchspacing NACH \selectlanguage: babel setzt den Satzabstand je
% Sprache, und british stellt den weiten englischen Abstand nach dem Punkt
% ein. Der faellt im Brief als Loecher auf und passt nicht zu den deutschen
% Dokumenten, die ngerman ohnehin auf frenchspacing stellt -- fuer Deutsch
% ist die Zeile daher wirkungslos.
\AtBeginDocument{\expandafter\selectlanguage\expandafter{\cvlanguage}%
    \cvhyphenexceptions\frenchspacing}

% Flatterrand statt Blocksatz: in schmalen Spalten erzeugt Blocksatz
% Trennungskaskaden und Löcher. \RaggedRight erlaubt Trennung noch,
% hält den Grauwert aber ruhig. \cvbodyhyphenpenalty kann die Inhaltsdatei
% überschreiben: 10000 schaltet die Trennung ab  --
% getrennte Wörter kamen in Bewerbungsportalen als "Haushaltsvoll- zugs-
% programm" an, weil die Portale den Text aus dem PDF auslesen.
\providecommand{\cvbodyhyphenpenalty}{200}
% Ohne Trennung dehnt \raggedright statt \RaggedRight: ragged2e lässt eine
% Zeile vor einem langen Wort lieber überlaufen ("Prozessdokumentation"
% ragte 13,6 pt in den Rand).
\newcommand{\bodyflow}{\RaggedRight\hyphenpenalty=\cvbodyhyphenpenalty\relax
    \ifnum\cvbodyhyphenpenalty<10000 \else\raggedright\fi}

%-------------------------------------------------------------------------
%  Icons
%-------------------------------------------------------------------------

% Die Symbolschrift hat keine Unicode-Zuordnung; ausgelesen wurde jedes
% Symbol zu einem Störzeichen vor der Adresse oder Telefonnummer ("Ȱ", "Л").
% \icotext gibt ihnen einen leeren Ersatztext -- Portale sehen nur die Angabe.
\newcommand{\icotext}[1]{\BeginAccSupp{method=plain,ActualText={}}#1\EndAccSupp{}}
\newcommand{\icoloc}{\icotext{\faMapMarker*}}
\newcommand{\icophone}{\icotext{\faPhone*}}
\newcommand{\icomail}{\icotext{\faEnvelope}}
\newcommand{\icolinked}{\icotext{\faLinkedin}}
\newcommand{\icocal}{\icotext{\faBirthdayCake}}
\newcommand{\icoheart}{\icotext{\faHeart}}
\newcommand{\icocar}{\icotext{\faCar}}

%-------------------------------------------------------------------------
%  Sechster Kontaktplatz — z.B. LinkedIn, GitHub oder der Führerschein.
%  Die Inhaltsdatei setzt Symbol, Text und optional eine URL; ist die URL
%  leer, steht der Text ohne Verlinkung da.
%-------------------------------------------------------------------------
\providecommand{\CVlinkicon}{\icolinked}
\providecommand{\CVlinktext}{}
\providecommand{\CVlinkurl}{}
\newcommand{\cvlink}{\ifdefempty{\CVlinkurl}{\CVlinktext}{\href{\CVlinkurl}{\CVlinktext}}}

%-------------------------------------------------------------------------
%  Haken für Abschnitte, die nur eine Person braucht (z.B. Referenzen).
%  Wird am Ende des Lebenslaufs ausgeführt; die Inhaltsdatei füllt ihn.
%-------------------------------------------------------------------------
\providecommand{\cvextrasections}{}

%-------------------------------------------------------------------------
%  Master-Lebenslauf und optionale Bausteine.
%
%  Die Inhaltsdatei hält ALLE Angaben. Was im Dossier erscheint, steuert
%  die meta.tex der Bewerbung:
%      \cvshow{schluessel}   zeigt einen Baustein, der sonst wegbleibt
%      \cvhide{schluessel}   lässt einen Baustein weg, der sonst erscheint
%  Die Schlüssel stehen im Kopf von themes/cv-content.tex. application.tex
%  und preview.tex definieren \cvshow und \cvhide, BEVOR sie die meta.tex
%  lesen; hier stehen sie nur als Rückfallebene für andere Treiber.
%
%  In der Inhaltsdatei:
%      \cvopt{schluessel}{...}        nur mit \cvshow oder im Master
%      \cvstd{schluessel}{...}        immer, ausser mit \cvhide (im Master immer)
%      \cvifopt{schluessel}{dann}{sonst}
%  \def\cvmaster{} (make cv-master über templates/cv-master.tex) zeigt
%  alles: Das ist der Master-Lebenslauf, aus dem je Stelle reduziert wird.
%-------------------------------------------------------------------------
\providecommand{\cvshow}[1]{\expandafter\def\csname cvshow@#1\endcsname{}}
\providecommand{\cvhide}[1]{\expandafter\def\csname cvhide@#1\endcsname{}}
\newcommand{\cvifopt}[3]{\ifdef{\cvmaster}{#2}{\ifcsdef{cvshow@#1}{#2}{#3}}}
\newcommand{\cvopt}[2]{\cvifopt{#1}{#2}{}}
\newcommand{\cvstd}[2]{\ifdef{\cvmaster}{#2}{\ifcsdef{cvhide@#1}{}{#2}}}

%-------------------------------------------------------------------------
%  Achter Kontaktplatz: GitHub (Schlüssel github in der
%  Inhaltsdatei, gedacht für Software- und Datenstellen). Leer gelassen
%  erscheint er in keinem Theme.
%-------------------------------------------------------------------------
\providecommand{\CVgithubtext}{}
\providecommand{\CVgithuburl}{}
\newcommand{\icogithub}{\icotext{\faGithub}}
\newcommand{\cvgithub}{\ifdefempty{\CVgithuburl}{\CVgithubtext}{\href{\CVgithuburl}{\CVgithubtext}}}

%-------------------------------------------------------------------------
%  Haken für einen Abschnitt direkt nach der Berufserfahrung (z.B. "Open Source & Projekte"). Leer = kein Abschnitt. Jedes
%  Theme ruft ihn in \makecvbody auf.
%-------------------------------------------------------------------------
\providecommand{\cvprojectsection}{}

%-------------------------------------------------------------------------
%  Zwischenüberschrift innerhalb einer Station (z.B. "Entwicklung" und
%  "Versuch & Prüfung"). Steht zwischen
%  zwei cvbullets-Listen und bleibt mit dem ersten Punkt darunter zusammen.
%  Vorgabe für alle Themes; kompakt setzt eine eigene Fassung.
%-------------------------------------------------------------------------
\providecommand{\cvsubhead}[1]{%
    \par\Needspace*{5\baselineskip}\addvspace{4pt}%
    {\nohyph\mediumfont #1}\par\nopagebreak
}

%-------------------------------------------------------------------------
%  Siebter Kontaktplatz — Nationalität und Bewilligung. In der Schweiz
%  gehört die Angabe bei ausländischen Bewerbenden in die Personalien
%  (Aufenthaltsbewilligung B/C, für Liechtenstein die Grenzgänger-
%  bewilligung). Leer gelassen erscheint die Zeile nicht. 
%-------------------------------------------------------------------------
\providecommand{\CVnationality}{}
\newcommand{\icoflag}{\icotext{\faPassport}}

%-------------------------------------------------------------------------
%  Neunter Kontaktplatz — Bewilligung als eigener Platz:
%  "Deutsch, Niederlassungsbewilligung C" in einem Platz lief im Kompakt-
%  Theme bis in den GitHub-Platz daneben. Mit \CVpermit stehen Nationalität
%  und Bewilligung getrennt. Leer gelassen bleibt alles wie bisher.
%-------------------------------------------------------------------------
\providecommand{\CVpermit}{}
\newcommand{\icopermit}{\icotext{\faIdCard}}

%-------------------------------------------------------------------------
%  Schrift des Brieftexts. Voreinstellung: die \normalsize des Themes,
%  also dieselbe Grösse wie der Lebenslauf. Die Inhaltsdatei kann den
%  Brief grösser setzen — bei Briefen von 200 bis 300 Wörtern bleibt die
%  Seite in 9,4 pt sonst zur Hälfte leer . Wird am Ende
%  von \makeletterhead ausgeführt, nach dem \parskip des Themes.
%-------------------------------------------------------------------------
\providecommand{\letterbodyfont}{}

%-------------------------------------------------------------------------
%  Unterschrift als Bild. Leer = keine. Die Briefe rufen \lettersignature
%  zwischen Grussformel und Namen auf; ohne Bild bleibt nur der Abstand.
%-------------------------------------------------------------------------
\providecommand{\CVsignature}{}
\newcommand{\lettersignature}{%
    \ifdefempty{\CVsignature}{\vspace{0.3cm}}{%
        \par\vspace{2pt}\noindent\includegraphics[height=0.95cm]{\CVsignature}\par\vspace{1pt}}}

%-------------------------------------------------------------------------
%  Bild-Helfer
%-------------------------------------------------------------------------
% Quadratischer Ausschnitt mit abgerundeten Ecken.
%   #1 Kantenlänge des Rahmens   #2 Eckenradius
% Zwei Stellschrauben kommen aus der Inhaltsdatei, weil sie zum Foto
% gehören und nicht zum Theme — beide relativ zur Kantenlänge, damit
% derselbe Wert in allen Themes stimmt, egal wie gross der Rahmen ist:
%   \CVphotozoom   Bildbreite als Vielfaches der Kantenlänge (1 = bündig)
%   \CVphotoshift  Motivversatz; negativ schiebt das Bild nach unten,
%                  das Gesicht rutscht also von oben in die Mitte
\providecommand{\CVphotozoom}{1}
\providecommand{\CVphotoshift}{0}
%   \CVphotomirror 1 = wie aufgenommen, -1 = seitenverkehrt (damit die Person zum
%                  Text hin blickt, wenn das Foto rechts steht)
\providecommand{\CVphotomirror}{1}
\newcommand{\photobox}[2]{%
    \begin{tikzpicture}
        \pgfmathsetlengthmacro{\pbhalf}{0.5*#1}
        \pgfmathsetlengthmacro{\pbwide}{\CVphotozoom*#1}
        \pgfmathsetlengthmacro{\pbshift}{\CVphotoshift*#1}
        \clip[rounded corners=#2] (-\pbhalf,-\pbhalf) rectangle (\pbhalf,\pbhalf);
        \node[anchor=center, inner sep=0] at (0,\pbshift)
            {\scalebox{\CVphotomirror}[1]{\includegraphics[width=\pbwide]{\CVphoto}}};
    \end{tikzpicture}%
}

%-------------------------------------------------------------------------
%  Zweiter Zeitraum derselben Rolle beim selben Arbeitgeber.
%
%  Wer eine Funktion zweimal innehatte, bekommt EINEN Eintrag mit zwei
%  Zeiträumen — nicht zwei Einträge mit identischer Aufgabenliste. Der
%  jüngere Zeitraum steht oben und bestimmt, wo der Eintrag im Lebenslauf
%  einsortiert wird; der ältere folgt als zweite Zeile in der Datumsspalte.
%
%      \cvrole{Immobilien-Bewirtschafterin}{Ing.~Bau~AG}{Vaduz}%
%             {2022 -- 2023\cvalso{2018 -- 2019}}
%
%  Bewusst ohne "und": die Datumsspalte des Akzent-Themes ist 2,2 cm breit,
%  "und 2018 -- 2019" bricht dort auf drei Zeilen um. Dass es zwei Einsätze
%  sind, sagt ohnehin die Zeile unter dem Rollentitel.
%-------------------------------------------------------------------------
\newcommand{\cvalso}[1]{\\ #1}

%-------------------------------------------------------------------------
%  Organisationszeile "Firma · Ort" — beide Teile optional, damit Einträge
%  ohne Arbeitgeber (Familienphase, Auszeit) keinen leeren Trenner zeigen.
%  Nutzt \orgfont, \metafont und die Farbe muted aus dem jeweiligen Theme.
%-------------------------------------------------------------------------
\newcommand{\cvorgline}[2]{%
    % Keine Silbentrennung: Firmen- und Ortsnamen gehören nicht zerlegt.
    % "Zwickau, Deutsch-|land" liest sich wie ein Setzfehler. Umbrüche an
    % den Wortzwischenräumen bleiben möglich.
    \nohyph
    \ifblank{#1}{%
        \ifblank{#2}{}{{\metafont\color{muted}#2}}%
    }{%
        {\orgfont#1}%
        \ifblank{#2}{}{\,{\metafont\color{muted}\textbullet}\,{\metafont\color{muted}#2}}%
    }%
}

%-------------------------------------------------------------------------
%  Links — klickbar, ohne Rahmen
%-------------------------------------------------------------------------
\usepackage[hidelinks, bookmarks=false, unicode]{hyperref}
\urlstyle{same}

%-------------------------------------------------------------------------
%  Weitere optionale Angaben. Leer gelassen erscheinen sie nirgends.
%    \CVkeywords      Stichwortzeile unter dem Titel (kompakt, rueckgrat)
%    \CVwebtext/-url  eigene Website als weiterer Link (rueckgrat)
%    \CVphotoaspect   Höhe des Fotos im Verhältnis zur Breite (rueckgrat),
%                     1 = quadratisch, 1.25 = Hochformat 4:5
%-------------------------------------------------------------------------
\providecommand{\CVkeywords}{}
\providecommand{\CVwebtext}{}
\providecommand{\CVweburl}{}
\newcommand{\cvweb}{\ifdefempty{\CVweburl}{\CVwebtext}{\href{\CVweburl}{\CVwebtext}}}
\providecommand{\CVphotoaspect}{1}

\providecommand{\LblProjects}{Projekte}
\providecommand{\LblReferences}{Referenzen}
\providecommand{\LblEnclosures}{Beilagen}
\providecommand{\LblContents}{Inhalt}
\providecommand{\LblApplicationAs}{Bewerbung als}
\providecommand{\LblCoverLetter}{Anschreiben}
\providecommand{\LblMotivation}{Motivationsschreiben}

%-------------------------------------------------------------------------
%  Beilagenzeile am Ende des Anschreibens. Erscheint nur, wenn die meta.tex
%  \newcommand{\enclosures}{Lebenslauf, Arbeitszeugnisse, Diplome} setzt.
%-------------------------------------------------------------------------
\newcommand{\letterenclosures}{\ifdefined\enclosures
    \par\vspace{0.5cm}\noindent{\metafont\color{muted}\LblEnclosures: \enclosures}\par\fi}

%-------------------------------------------------------------------------
%  Deckblatt (\cvshow{coverpage} in der meta.tex): Briefkopf des Themes,
%  die Zielstelle in der Seitenmitte, unten das Inhaltsverzeichnis des
%  Dossiers. \coverpagecontents setzt application.tex je nach Schaltern.
%-------------------------------------------------------------------------
\providecommand{\coverpagecontents}{\LblCV}
% Grösse der Zielstelle auf dem Deckblatt, vom Theme überschreibbar
\providecommand{\coverpagetitlesize}{\fontsize{24pt}{29pt}\selectfont}
\newcommand{\makecoverpage}{%
    \makeletterhead
    \setlength{\parskip}{0pt}%
    \vspace*{\fill}%
    \noindent{\metafont\color{muted}\MakeUppercase{\LblApplicationAs}}\par\vspace{8pt}%
    \noindent{\nohyph\displayfont\coverpagetitlesize\color{headercolor}\positionname\par}%
    \vspace{6pt}%
    \noindent{\fontsize{13pt}{17pt}\selectfont\companyname\par}%
    \vspace*{\fill}\vspace*{\fill}%
    \noindent{\metafont\color{muted}\MakeUppercase{\LblContents}}\par\vspace{4pt}%
    \noindent{\color{rule}\rule{\linewidth}{0.5pt}}\par\vspace{6pt}%
    \noindent\coverpagecontents\par
    \newpage
}

\usepackage{eso-pic}

% Flatterrand mit wenig Trennungen: In einer Schrift mit fester Zeichenbreite fallen
% Trennstriche stärker auf als kurze Zeilen.
\renewcommand{\bodyflow}{\RaggedRight\hyphenpenalty=1000}

%-------------------------------------------------------------------------
%  Farben — Grün nur für Knoten, nie für Text
%-------------------------------------------------------------------------

\definecolor{ink}{HTML}{1A1A17}      % Name, Titel der Einträge
\definecolor{inksoft}{HTML}{4A4A44}  % Fliesstext
\definecolor{muted}{HTML}{67665F}    % Daten, Schlüssel, Kontakt
\definecolor{rule}{HTML}{D9D8D1}     % Haarlinien
\definecolor{accent}{HTML}{009E73}   % Knoten auf dem Rückgrat

% Kompatibilität zu den bestehenden Brief-Dateien
\colorlet{headercolor}{ink}
\colorlet{linecolor}{ink}

%-------------------------------------------------------------------------
%  Schrift — JetBrains Mono NL (SIL Open Font License, fonts/OFL-JetBrainsMono.txt)
%-------------------------------------------------------------------------

\setmainfont{JetBrainsMonoNL}[
    Path = fonts/, Extension = .ttf,
    UprightFont    = *-Regular,
    ItalicFont     = *-Italic,
    BoldFont       = *-SemiBold,
    BoldItalicFont = *-BoldItalic]
% \sffamily muss dieselbe Schrift liefern — die Briefdateien verwenden es
\setsansfont{JetBrainsMonoNL}[
    Path = fonts/, Extension = .ttf,
    UprightFont = *-Regular, ItalicFont = *-Italic, BoldFont = *-SemiBold]
\newfontfamily\displayfont{JetBrainsMonoNL}[
    Path = fonts/, Extension = .ttf, UprightFont = *-SemiBold]
\newfontfamily\mediumfont{JetBrainsMonoNL}[
    Path = fonts/, Extension = .ttf, UprightFont = *-Medium, ItalicFont = *-MediumItalic]

% Grundschriftgrad des Dokuments
\makeatletter
\renewcommand\normalsize{\@setfontsize\normalsize{9.3}{12.9}}
\makeatother
\AtBeginDocument{\normalsize\color{inksoft}}

\newcommand{\namefont}{\mediumfont\fontsize{25pt}{29pt}\selectfont}
\newcommand{\rolelabelfont}{\fontsize{11pt}{15pt}\selectfont}
\newcommand{\metafont}{\fontsize{7.8pt}{11pt}\selectfont}
\newcommand{\datefont}{\fontsize{7.8pt}{11pt}\selectfont}
\newcommand{\rolefont}{\displayfont\fontsize{10pt}{13.2pt}\selectfont}
\newcommand{\orgfont}{\fontsize{9.3pt}{12.9pt}\selectfont\color{muted}}
\newcommand{\notefont}{\fontsize{8pt}{11pt}\selectfont\itshape}
\newcommand{\labelfont}{\mediumfont\fontsize{8.6pt}{11pt}\selectfont
    \addfontfeature{LetterSpace=2.5}}
\newcommand{\lead}[1]{{\mediumfont#1}}

% Grundlinien-Schutz, siehe kompakt.tex
\newcommand{\vtopsafe}{\leavevmode}

%-------------------------------------------------------------------------
%  Linke Spalte, Knoten, Rückgrat
%-------------------------------------------------------------------------

% Inhalt der linken Spalte am Anfang der laufenden Zeile, rechtsbündig. Die Box
% bekommt Höhe und Tiefe null, damit sie die Zeilenhöhe nie verändert.
\newsavebox{\rglcbox}
\newcommand{\rgleftcol}[1]{\leavevmode\sbox{\rglcbox}{\parbox[t]{\rgcolL}{%
    \RaggedLeft\leavevmode #1}}%
    \ht\rglcbox=0pt \dp\rglcbox=0pt \llap{\usebox{\rglcbox}\hspace{\rgcolG}}}

% Grüner Knoten auf dem Rückgrat, mit weissem Hof: Die Linie scheint dahinter durchzulaufen
\newcommand{\rgnode}{\rlap{\hspace{-0.5\rgcolG}\raisebox{0.3em}{\tikz[overlay]{%
    \fill[white] (0,0) circle (4.3pt); \fill[accent] (0,0) circle (2.9pt);}}}}

% Kleiner grauer Punkt als Aufzählungszeichen
\newcommand{\rgdot}{\tikz[baseline=-0.5ex]\fill[muted] (0,0) circle (1.1pt);}

% Das Rückgrat: eine senkrechte Linie durch jede Seite, hinter den Knoten
\AddToShipoutPictureBG{\AtPageLowerLeft{%
    \put(\LenToUnit{\dimexpr\rgmarginleft+\rgcolL+0.5\rgcolG-0.4pt\relax},\LenToUnit{9.5mm}){%
        \color{ink}\rule{0.8pt}{\dimexpr\paperheight-9.5mm-\ifnum\value{page}>1 8.5mm\else 18mm\fi\relax}}}}

%-------------------------------------------------------------------------
%  Sektionen — Titel in der linken Spalte, Knoten auf der Linie
%-------------------------------------------------------------------------

\newcommand{\rglabel}[1]{\rgleftcol{\nohyph\labelfont\color{ink}#1}\rgnode}

% Vor Einträgen und Fakten-Zeilen steht der Titel auf einer eigenen Zeile.
\newcommand{\cvsection}[1]{%
    \par\addvspace{14pt}%
    \Needspace*{5\baselineskip}%
    {\rglabel{#1}\strut\par}\nobreak\vspace{5pt}\nobreak
}

%-------------------------------------------------------------------------
%  Einträge — Datum in der linken Spalte, Inhalt rechts.
%  Die Einträge sind gewöhnliche Absätze und laufen über Seitenumbrüche.
%-------------------------------------------------------------------------

\newenvironment{cvexperience}{}{\par}

\newcommand{\cvrole}[4]{%
    \par\addvspace{7pt}%
    \Needspace*{4\baselineskip}% Station nicht als Waise an den Seitenfuss
    {\rgleftcol{\datefont\color{muted}#4}\nohyph\rolefont\color{ink}#1\par}\nobreak
    {\leavevmode\cvorgline{#2}{#3}\par}\nobreak
    \addvspace{1.5pt}%
}

\newcommand{\cvnote}[1]{{\bodyflow\notefont\color{muted}#1\par}\vspace{1pt}}

\newcommand{\cvplain}[3]{%
    \par\addvspace{7pt}%
    \Needspace*{3\baselineskip}%
    {\rgleftcol{\datefont\color{muted}#2}\nohyph\rolefont\color{ink}#1\par}\nobreak
    {\bodyflow #3\par}%
}

\newenvironment{cvbullets}{%
    \begin{itemize}[label=\rgdot, leftmargin=10pt, labelsep=5pt,
                    itemsep=1.8pt, parsep=0pt, topsep=3pt, partopsep=0pt]
        \bodyflow
    }{\end{itemize}}

%-------------------------------------------------------------------------
%  Fakten-Zeilen (Kompetenzen, Sprachen) — Schlüssel klein in der linken Spalte
%-------------------------------------------------------------------------

\newcommand{\cvfact}[2]{%
    \par{\rgleftcol{\metafont\color{muted}\MakeLowercase{#1}}\bodyflow #2\par}%
    \addvspace{2.2pt}%
}

\newcommand{\cvtag}[1]{#1\quad}  % Chips braucht dieses Theme nicht

%-------------------------------------------------------------------------
%  Kopf des Lebenslaufs
%-------------------------------------------------------------------------

% Adresse am ersten Komma umbrechen: "Strasse 1, 8000 Ort" wird zweizeilig.
\makeatletter
\def\rg@split#1, #2\@nil{#1\ifx\relax#2\relax\else\newline\rg@strip#2\@nil\fi}
\def\rg@strip#1, \@nil{#1}
\newcommand{\rgaddress}{\expandafter\rg@split\CVaddress, \@nil}
\makeatother

% "geb." vor dem Geburtsdatum, englisch "born"
\newcommand{\rgborn}{\ifdefstring{\cvlanguage}{ngerman}{geb.\ }{born\ }}

% Reihenfolge im Textstrom (für Parser): Name, Titel, Link, Foto, dann die Kontaktspalte.
% Die Kontaktspalte wird danach nach links zurückgeschoben.
\newcommand{\makecvheader}{%
    \noindent
    \begin{minipage}[t]{\dimexpr\linewidth-\rgphotow-\rgcolG\relax}\vspace{0pt}%
        \vspace{-3pt}{\namefont\nohyph\color{ink}\CVname\par}\vspace{2pt}%
        {\rolelabelfont\color{inksoft}\CVrole\par}%
        \vspace{39pt}%
        {\metafont\color{muted}\cvlink\par}%
    \end{minipage}%
    \hspace{\rgcolG}%
    \begin{minipage}[t]{\rgphotow}\vspace{0pt}%
        \photobox{\rgphotow}{0pt}%
    \end{minipage}%
    \hspace*{-\dimexpr\linewidth+\rgpull\relax}%
    \begin{minipage}[t]{\rgcolL}\vspace{0pt}\RaggedLeft
        \vspace{28.6pt}% erste Kontaktzeile auf der Grundlinie des Titels
        {\metafont\color{muted}%
         \rgaddress\par
         \rgborn\CVbirth\par
         \CVcivil\par
         \vspace{6pt}%
         \CVphone\par
         \href{mailto:\CVmail}{\CVmail}\par}%
    \end{minipage}%
    \par\vspace{12pt}%
    % Profil: Der Titel teilt die Zeile mit dem Text
    \Needspace*{4\baselineskip}%
    {\rglabel{\LblProfile}\bodyflow\ContentSummary\par}%
}

%-------------------------------------------------------------------------
%  Seitenstil — ab Seite 2 Name und Dokument im Kopf, Seitenzahl links unten
%-------------------------------------------------------------------------

\fancypagestyle{cvstyle}{%
    \fancyhf{}
    \fancyhfoffset[L]{\rgpull}
    \renewcommand{\headrulewidth}{0pt}
    \renewcommand{\footrulewidth}{0pt}
    \fancyhead[L]{\metafont\color{muted}\makebox[\rgcolL][r]{\CVname}\hspace{\rgcolG}\LblCV}
    \fancyfoot[L]{\metafont\color{muted}\makebox[\rgcolL][r]{\thepage}}
}
\fancypagestyle{cvfirst}{%
    \fancyhf{}
    \fancyhfoffset[L]{\rgpull}
    \renewcommand{\headrulewidth}{0pt}
    \renewcommand{\footrulewidth}{0pt}
}

%-------------------------------------------------------------------------
%  Briefkopf (Anschreiben / Motivationsschreiben) — wie der Kopf des
%  Lebenslaufs, ohne Foto und ohne persönliche Angaben
%-------------------------------------------------------------------------

\newcommand{\makeletterhead}{%
    \pagestyle{cvfirst}\thispagestyle{cvfirst}%
    \setlength{\parskip}{7.5pt}% Brieftext: Absätze durch Weissraum getrennt
    \noindent
    \begin{minipage}[t]{\linewidth}\vspace{0pt}\setlength{\parskip}{0pt}%
        \vspace{-3pt}{\namefont\nohyph\color{ink}\CVname\par}\vspace{2pt}%
        {\rolelabelfont\color{inksoft}\CVrole\par}%
    \end{minipage}%
    \hspace*{-\dimexpr\linewidth+\rgpull\relax}%
    \begin{minipage}[t]{\rgcolL}\vspace{0pt}\RaggedLeft\setlength{\parskip}{0pt}%
        \vspace{28.6pt}%
        {\metafont\color{muted}%
         \rgaddress\par
         \CVphone\par
         \href{mailto:\CVmail}{\CVmail}\par}%
    \end{minipage}%
    \par\vspace{14mm}%
    \bodyflow % Brieftext im Flatterrand: Blocksatz reisst in dieser Schrift Löcher
}

% Zwischentitel im Brief: steht in der linken Spalte neben dem nächsten Absatz
\newcommand{\lettersection}[1]{%
    \par\addvspace{11pt}%
    \everypar{\rglabel{#1}\everypar{}}%
}

%-------------------------------------------------------------------------
%  Rumpf des Lebenslaufs
%-------------------------------------------------------------------------

\newcommand{\makecvbody}{%
    \pagestyle{cvstyle}\thispagestyle{cvfirst}
    \makecvheader

    \cvsection{\LblSkills}
    \ContentSkills

    \cvsection{\LblExperience}
    \ContentExperience

    \cvsection{\LblEducation}
    \ContentEducation

    \cvsection{\LblFurther}
    \ContentFurther

    \cvsection{\LblLanguages}
    \ContentLanguages

    \cvextrasections
}
